\subsection{二つの葉付き支持の排除から再利用する方法}
\label{nl19:section}

$G_{5,7}$ と $G_{5,9}$ は，一高完全標準形，三高 matching，
二高の正負二分，共通低 puncture の signed-divisor 分類，node 非負性と対向 $D$ 相殺で排除された．
標準 $K_5$ graph，四直線 cones，marked punctures を保つ五次の整数接着問題に関する結論である．
葉を $4$，その高隣接を $0$ とし，葉の標準 rays を $(0,B,C,D)$ とする．
$Q=th(h+1)+1$，$u=th$，$v=t(h+1)$，$cd=0$ の全整数自由度を保持する．

$G_{5,7}$ では $0$ は三高で，その唯一の低 ray の位置 $m$ が $B,C,D$ の全配置を尽くす．
可能な配置は
\[
 (B,N_{0B}=2,Q\le5),\quad(C,N_{0C}=1,u\le4),\quad
 (C,N_{0C}=2,h=t=1),\quad(D,N_{0D}=1,v\le4)
\]
に限られ，従って $Q\le21$ である．
$G_{5,9}$ では $0$ は四高で，残る三高 core の位置 $J=B,C,D$ を全て残す．
$B$ の特別な reverse は $J=B$ なら $1,2$，他なら $1,2,-2$ だけとなる．
$-2$ の枝は $Q=3$ に限られ，三配置の低低 entry と matching の比較から $Q\le13$ を得る．
いずれの上界も，当該 support の仮定を用いて証明した結果であり，数値探索の上限ではない．

再利用する葉接着の式は
\[
 \ell=(Q,1,1,1),\quad R_{ij}=p^i_{4j},\quad X_{ij}=p^i_{j4},\quad V=P^4,
\]
\[
 R_{ij}=\ell_i(1-N_{ij})-V_{ij},\quad\sum_{j\ne i}N_{ij}=5,\quad
 R_{ij}X_{ij}=\ell_iq_{ij}-1,
\]
\[
 \sum_{j\ne i}X_{ij}+\sum_{j\ne i}V_{ji}=-2\ell_i.
\]
高 ray を含む葉の pair product は $Q-1$ であり，零ではない．
各 node 行は二十一通りだが，それだけでは自由整数も他の分母も有限にならない．
固定 reverse の相手は一点または全偶奇半直線で処理し，中央辺では
\[
 (Qr+R-\kappa Rr)q=R+r
\]
を用いる．係数と右辺がともに零なら，
$q=1+\operatorname{lcm}(2,|r|)a$，$a\ge1$ の全族を残す．

$G_{5,7}$ の残余は $160$ partner 型，中央辺の選択まで含めて $212$ 型である．
中心の $D$ 式に必要な情報をまとめた十三行では，自由分母の七行が $q=1,-1,-3$ を要求し，
残る六行は非零の $D$ 残差を持つ．
$G_{5,9}$ では点 $908$ 型と全半直線 $18$ 型が残る．後者は二高 core の canonical な
第二高分母 $b=3$ を要求して矛盾する．低低 pair の逆 entry は別の低辺の node 予算から
$0,-1,-3,-4$ に限られ，点型から $2596$ 個の $D$ 線形系を得る．
$2563$ 個は有理係数で不整合，残る $33$ 個は $3a+3b=-6$，$a,b\ge1$ を要求する．
各矛盾の元の四式による有理線形結合を証明書に保存した．
自由分母が多数残っても，先に全局所行列を完成させず，$D$ の必要条件だけで排除できることを示す．

\paragraph{小さい正 entry と，局所例の使い分け．}
二高 core の二つの低行の正 entries が同時に二以下になることはない．
canonical には $r,T$ であり，正枝は $T>r+1$，負枝の唯一の小候補は
$(s-4)(z+3)=-6$ と整除性に反する．
この補題は，隣接 matching core から node によって小さい上界が伝わる場合に使える．
一方，三高 core の低行 $(-1,-1,1)$ 自体は不可能ではない．例えば
\[
 q=(1,3,31,7),\qquad
 P=\begin{pmatrix}
 0&-1&-1&1\\-2&0&4&-5\\-30&23&0&-24\\6&-4&-9&0
 \end{pmatrix}
\]
は row/pair equations と正値性を満たす．正二閉路の gains は $6/7$ と $92/93$ である．
葉型の排除には隣接 core と node/$D$ を用いることを忘れない．

完全証明は算術稿の \nolinkurl{g5_7_obstruction_v19.tex} と
\nolinkurl{g5_9_obstruction_v19.tex}，型の全証明書は同名の主検算と結果 JSON にある．
型数は canonical branches の重複を残す手順上の個数であり，同型類数ではない．
